\chapter{Workbook Phần 6: 20 Bài Tập Thực Hành Data Engineering, Big Data \& MLOps}

\begin{lythuyetbox}[Mục Tiêu Luyện Tập Phần 6]
Phần này cung cấp \textbf{20 bài tập thực hành kỹ thuật chuyên sâu} dành cho kỹ sư dữ liệu (Data Engineer) và kỹ sư MLOps. Nội dung bao quát toàn diện từ mô hình hóa dữ liệu kho Kimball (Star Schema, SCD2), thiết kế đường ống ETL/ELT có tính lũy đẳng (Idempotency), Docker đa tầng, Apache Airflow TaskFlow, dbt models/tests, tối ưu hóa PySpark với Broadcast Join, Delta Lake ACID, kiến trúc luồng dữ liệu Kafka, triển khai Microservice FastAPI độ trễ thấp, và tự động hóa CI/CD với GitHub Actions.
\end{lythuyetbox}

\begin{thuchanhbox}[Tệp Dữ Liệu \& Kịch Bản Thực Hành DE \& MLOps Trực Tiếp (Bấm Vào Để Mở)]
Bạn có thể trực tiếp thực thi, kiểm thử và mở mã nguồn các bài tập DE và MLOps qua các tệp đã xây dựng sẵn:
\begin{itemize}
    \item \href{run:./data/raw_transactions.csv}{\Large\textbf{\color{navyprimary}data/raw\_transactions.csv}} -- 10,000 giao dịch thô chứa bản ghi lỗi: \texttt{INVALID\_STATUS}, số tiền âm, mã giảm giá sai lệch.
    \item \href{run:./data/analytics_warehouse.db}{\Large\textbf{\color{navyprimary}data/analytics\_warehouse.db}} -- Cơ sở dữ liệu Data Warehouse SQLite chứa bảng hạt nhân \texttt{fact\_transactions} (9,659 dòng đã làm sạch).
    \item \href{run:./code/ch07_etl/batch_etl.py}{\Large\textbf{\color{navyprimary}code/ch07\_etl/batch\_etl.py}} -- Pipeline Batch ETL sản phẩm: Trích xuất, kiểm định chất lượng, lọc bỏ 341 lỗi, nạp kho lũy đẳng.
    \item \href{run:./code/ch08_orchestration/ecommerce_dag.py}{\Large\textbf{\color{navyprimary}code/ch08\_orchestration/ecommerce\_dag.py}} -- Apache Airflow DAG điều phối luồng xử lý tự động với TaskFlow API.
    \item \href{run:./code/ch09_ml/train_churn.py}{\Large\textbf{\color{navyprimary}code/ch09\_ml/train\_churn.py}} -- Pipeline huấn luyện Random Forest dự đoán Churn, đánh giá ROC-AUC và xuất artifact.
    \item \href{run:./code/ch12_mlops/app.py}{\Large\textbf{\color{navyprimary}code/ch12\_mlops/app.py}} -- Microservice FastAPI độ trễ thấp phục vụ suy luận thời gian thực cho mô hình ML.
\end{itemize}
\small\textit{(Xem trực tuyến tại: \url{https://github.com/TontonYuta/DataLearning/tree/main/code})}
\end{thuchanhbox}

\begin{baitapbox}[Hệ Thống 20 Bài Tập Thực Chiến DE \& MLOps]
\begin{enumerate}
    \item \textbf{Bài 1 (Thiết Kế Kimball Star Schema)}: Cho hệ thống E-commerce với hàng triệu giao dịch. Hãy thiết kế lược đồ hình sao gồm: Bảng Fact \texttt{fact\_order\_items} (hạt nhân Grain: từng dòng sản phẩm trong đơn) và 4 bảng Dimension (\texttt{dim\_customer}, \texttt{dim\_product}, \texttt{dim\_date}, \texttt{dim\_store}). Chỉ ra khóa thay thế (Surrogate Key) và một Degenerate Dimension trong bảng Fact.
    \item \textbf{Bài 2 (Xử Lý Lịch Sử Biến Đổi Chậm SCD Type 2)}: Khi khách hàng chuyển nhà từ Hà Nội vào TP.HCM, làm thế nào để lưu vết lịch sử nhằm đảm bảo các báo cáo doanh thu theo vùng trong quá khứ không bị sai lệch? Viết cấu trúc DDL của bảng \texttt{dim\_customer} có chứa các cột kỹ thuật quản lý SCD2 (\texttt{valid\_from}, \texttt{valid\_to}, \texttt{is\_current}).
    \item \textbf{Bài 3 (Nguyên Tắc Lũy Đẳng \& Làm Sạch Dữ Liệu Trên \texttt{batch\_etl.py})}: Mở script \href{run:./code/ch07_etl/batch_etl.py}{\texttt{code/ch07\_etl/batch\_etl.py}} và tệp nguồn \href{run:./data/raw_transactions.csv}{\texttt{data/raw\_transactions.csv}} (10,000 dòng). Trình bày cơ chế lọc bỏ các bản ghi không đạt chuẩn (ví dụ trạng thái \texttt{'INVALID\_STATUS'}, chiết khấu âm). \textit{Kết quả kiểm tra}: Pipeline loại bỏ chính xác \textbf{341 bản ghi lỗi} và nạp \textbf{9,659 bản ghi sạch} vào kho \href{run:./data/analytics_warehouse.db}{\texttt{analytics\_warehouse.db}}.
    \item \textbf{Bài 4 (Xác Thực Schema Dữ Liệu Với Pydantic)}: Viết một schema Pydantic bằng Python để xác thực luồng dữ liệu thanh toán từ Webhook bên ngoài. Yêu cầu: \texttt{transaction\_id} kiểu chuỗi không rỗng, \texttt{amount} kiểu float dương ($> 0$), \texttt{currency} chỉ nhận giá trị trong tập \texttt{["VND", "USD"]}, và \texttt{created\_at} phải ở định dạng chuẩn ISO datetime.
    \item \textbf{Bài 5 (Tối Ưu Hóa Dockerfile Đa Tầng -- Multi-Stage Build)}: Một ứng dụng Ingestion viết bằng Python có kích thước Docker Image lên tới 1.2GB khi build với image gốc \texttt{ubuntu:latest} và gài đầy đủ trình biên dịch C/C++. Hãy viết Dockerfile tối ưu sử dụng \texttt{python:3.11-slim}, kỹ thuật Multi-stage build, và cấu hình chạy bằng người dùng không có quyền root (\texttt{non-root user}) để nén kích thước xuống dưới 150MB.
    \item \textbf{Bài 6 (Thiết Kế Cụm Dịch Vụ Với Docker Compose)}: Viết file cấu hình \texttt{docker-compose.yml} để khởi chạy đồng bộ 3 dịch vụ: PostgreSQL Database (lưu dữ liệu có mount volume bền vững), Redis Cache (cổng 6379), và một container Python Worker. Yêu cầu cấu hình cơ chế kiểm tra sức khỏe (\texttt{healthcheck}) đảm bảo Worker chỉ khởi động khi PostgreSQL đã sẵn sàng nhận kết nối.
    \item \textbf{Bài 7 (Apache Airflow TaskFlow API Trên \texttt{ecommerce\_dag.py})}: Mở kịch bản \href{run:./code/ch08_orchestration/ecommerce_dag.py}{\texttt{code/ch08\_orchestration/ecommerce\_dag.py}}. Viết một DAG Airflow hoàn chỉnh bằng cú pháp hiện đại TaskFlow (\texttt{@dag} và \texttt{@task}) thực hiện quy trình 3 bước: \texttt{extract()} (trích xuất đơn hàng), \texttt{transform(raw\_data)} (lọc đơn hủy và tính thuế), \texttt{load(clean\_data)} (ghi vào Database). Cấu hình tự động retry 3 lần.
    \item \textbf{Bài 8 (Airflow Trigger Rule \& Cảnh Báo Lỗi)}: Trong một pipeline phức tạp có nhiều nhánh rẽ, bạn muốn có một task cuối cùng gửi thông báo tổng kết lên kênh Slack bất kể các task phía trước thành công hay thất bại. Hãy cho biết cần thiết lập tham số \texttt{trigger\_rule} của task này là gì và giải thích cơ chế hoạt động.
    \item \textbf{Bài 9 (Xây Dựng dbt Staging Model)}: Viết mô hình dbt SQL \texttt{stg\_orders.sql} đọc dữ liệu thô từ bảng nguồn \texttt{raw.jaffle\_shop.orders}. Thực hiện các thao tác: đổi tên cột kiểu camelCase sang snake\_case, ép kiểu \texttt{order\_date} sang kiểu DATE, chuyển đổi giá tiền từ cents sang dollars, và lọc bỏ các đơn hàng thử nghiệm có \texttt{is\_test = true}.
    \item \textbf{Bài 10 (Kiểm Thử Chất Lượng Dữ Liệu Bằng dbt Tests)}: Viết file cấu hình \texttt{schema.yml} trong dbt để khai báo các bài test tự động cho bảng \texttt{stg\_customers}: Cột \texttt{customer\_id} phải là duy nhất (\texttt{unique}) và không được rỗng (\texttt{not\_null}); Cột \texttt{status} chỉ được phép nhận các giá trị định trước (\texttt{accepted\_values}). Viết thêm 1 bài Singular Test SQL kiểm tra không có đơn hàng nào có tổng tiền âm ($< 0$).
    \item \textbf{Bài 11 (Chiến Lược Cập Nhật Tăng Dần -- dbt Incremental)}: Khi bảng đơn hàng tích lũy hơn 100 triệu dòng, việc build lại toàn bộ (full-refresh) mỗi ngày sẽ gây cạn kiệt tài nguyên máy chủ. Hãy cấu hình cấu trúc model dbt \texttt{materialized='incremental'} với khóa duy nhất \texttt{unique\_key='order\_id'} và viết đoạn code Jinja lọc lấy các bản ghi có \texttt{updated\_at} mới hơn giá trị lớn nhất hiện có trong bảng đích.
    \item \textbf{Bài 12 (PySpark Schema Enforcement \& Partitioning)}: Viết đoạn mã PySpark đọc một tệp CSV lớn (dung lượng 50GB) bằng cách định nghĩa trước cấu trúc Schema tường minh (\texttt{StructType}) thay vì dùng \texttt{inferSchema=True}. Sau đó, thực hiện ghi dữ liệu ra định dạng Parquet nén Snappy, phân vùng (partition) theo 2 cột \texttt{year} và \texttt{month}.
    \item \textbf{Bài 13 (Tối Ưu Hóa Phép Join Bằng Broadcast Hash Join)}: Trong Spark, khi thực hiện phép Join giữa bảng \texttt{fact\_transactions} (1 tỷ bản ghi, kích thước 500GB) và bảng \texttt{dim\_store} (500 cửa hàng, kích thước 2MB), hiện tượng xáo trộn dữ liệu (Data Shuffling) qua mạng sẽ làm chậm pipeline nghiêm trọng. Viết câu lệnh Spark tối ưu bằng \texttt{broadcast()} và phân tích nguyên lý hoạt động của Broadcast Hash Join.
    \item \textbf{Bài 14 (PySpark Window Functions)}: Cho DataFrame \texttt{df\_sales} gồm \texttt{[department, employee\_id, salary, sale\_amount]}. Hãy viết mã PySpark sử dụng Window Specification để: (1) Tính tổng lũy kế doanh số (Running Total) của từng nhân viên theo thời gian; (2) Xếp hạng top 3 nhân viên có doanh số cao nhất trong từng phòng ban bằng hàm \texttt{dense\_rank()}.
    \item \textbf{Bài 15 (Kiến Trúc Lakehouse Với Delta Lake ACID)}: Trình bày sự khác biệt cốt lõi giữa Parquet thuần túy và Delta Lake. Viết câu lệnh SQL/PySpark để thực hiện tính năng \textit{Time Travel} (truy vấn lại dữ liệu tại phiên bản quá khứ \texttt{VERSION AS OF 3} hoặc tại thời điểm ngày hôm qua) và giải thích tác dụng của lệnh \texttt{VACUUM}.
    \item \textbf{Bài 16 (Hạ Tầng Luồng Dữ Liệu Thời Gian Thực Với Apache Kafka)}: Một hệ thống clickstream người dùng tạo ra 10,000 sự kiện/giây. Để đảm bảo thứ tự thời gian của các sự kiện do cùng một người dùng tạo ra (in-order delivery) khi ghi vào Kafka Topic có 12 partitions, bạn cần thiết lập tham số gì khi Producer gửi message? Hiện tượng Consumer Lag là gì và cách mở rộng quy mô (Scale out) nhóm Consumer?
    \item \textbf{Bài 17 (Triển Khai Microservice FastAPI Trên \texttt{code/ch12\_mlops/app.py})}: Mở tệp \href{run:./code/ch12_mlops/app.py}{\texttt{code/ch12\_mlops/app.py}}. Viết mã nguồn FastAPI nạp mô hình \texttt{churn\_model.joblib} qua cơ chế \texttt{lifespan context manager}, tiếp nhận payload xác thực Pydantic và trả về xác suất rời bỏ kèm nhãn phân loại (0 hoặc 1).
    \item \textbf{Bài 18 (Kiểm Thử Tải Hệ Thống Phục Vụ ML Bằng Locust)}: Viết một script kiểm thử tải (Load Testing) bằng công cụ Locust trong Python để giả lập 200 người dùng đồng thời gửi request liên tục đến endpoint \texttt{/predict} của mô hình ML. Đo đạc các chỉ số: Requests Per Second (RPS), Thời gian phản hồi trung vị (Median Response Time), và phân vị P99.
    \item \textbf{Bài 19 (Thiết Lập Đường Ống CI/CD Tự Động Với GitHub Actions)}: Viết file workflow cấu hình GitHub Actions \texttt{.github/workflows/data\_pipeline\_ci.yml} tự động kích hoạt khi có Pull Request vào nhánh \texttt{main}. Quy trình gồm: Khởi tạo môi trường Ubuntu, cài đặt Python 3.11, kiểm tra chất lượng mã nguồn bằng \texttt{flake8}, và tự động chạy toàn bộ bài kiểm thử đơn vị với \texttt{pytest}.
    \item \textbf{Bài 20 (Giám Sát Độ Lệch Dữ Liệu MLOps -- Data Drift Detection)}: Sau khi triển khai mô hình Machine Learning lên môi trường Production được 3 tháng, độ chính xác của mô hình bị suy giảm nghiêm trọng. Giải thích khái niệm Data Drift và Concept Drift. Trình bày phương pháp định lượng độ lệch dữ liệu đầu vào sử dụng chỉ số PSI (Population Stability Index) và ngưỡng kích hoạt huấn luyện lại mô hình (Retraining Trigger).
\end{enumerate}
\end{baitapbox}

\begin{dapanbox}[Đáp Án \& Lời Giải Ngắn Gọn Cho 20 Bài Tập DE \& MLOps]
\begin{itemize}
    \item \textbf{Bài 1}: Bảng Fact: \texttt{fact\_order\_items(order\_item\_key, order\_number, customer\_key, product\_key, date\_key, store\_key, quantity, unit\_price, discount\_amount, net\_amount)}. Surrogate Key: Các khóa nguyên tự tăng \texttt{*\_key}. Degenerate Dimension: \texttt{order\_number} (mã hóa đơn lưu trực tiếp trên Fact mà không cần tạo bảng Dim riêng vì không có thuộc tính mô tả bổ sung).
    \item \textbf{Bài 2}: DDL bảng SCD2: \texttt{CREATE TABLE dim\_customer (customer\_sk BIGINT PRIMARY KEY, customer\_id INT, full\_name VARCHAR(100), city VARCHAR(50), valid\_from DATE NOT NULL, valid\_to DATE, is\_current BOOLEAN NOT NULL)}. Khi khách đổi địa chỉ: UPDATE bản ghi cũ gán \texttt{valid\_to = CURRENT\_DATE - 1, is\_current = FALSE}; INSERT bản ghi mới với địa chỉ mới, \texttt{valid\_from = CURRENT\_DATE, valid\_to = '9999-12-31', is\_current = TRUE}.
    \item \textbf{Bài 3}: Áp dụng mô hình \textbf{Idempotent Ingestion}: Trong \texttt{batch\_etl.py}, 10,000 dòng được đọc; 341 dòng lỗi bị loại bỏ bằng bộ lọc \lstinline|df['order_status'].isin(['COMPLETED', 'PENDING', 'CANCELLED']) & (df['amount'] > 0) & (df['discount'] >= 0)|; 9,659 bản ghi hợp lệ được nạp vào kho \texttt{analytics\_warehouse.db} bằng cơ chế ghi đè phân vùng (\texttt{if\_exists='replace'} hoặc \texttt{INSERT OR REPLACE INTO}). Dù Airflow chạy lại nhiều lần thì kho dữ liệu luôn nhất quán.
    \item \textbf{Bài 4}: Khai báo Pydantic Schema:
    \begin{lstlisting}[language=python]
from pydantic import BaseModel, Field
from datetime import datetime
from typing import Literal

class TransactionPayload(BaseModel):
    transaction_id: str = Field(..., min_length=1)
    amount: float = Field(..., gt=0)
    currency: Literal["VND", "USD"]
    created_at: datetime
    \end{lstlisting}
    \item \textbf{Bài 5}: Sử dụng \texttt{python:3.11-slim}, Stage 1 (builder) cài đặt wheel packages; Stage 2 (final) chỉ sao chép wheel đã build, dọn dẹp cache \texttt{--no-cache-dir}, tạo người dùng \texttt{RUN useradd -m appuser \&\& USER appuser}. Image cuối cùng không chứa gcc hay header files, giảm kích thước từ 1.2GB xuống $\sim 115$MB.
    \item \textbf{Bài 6}: Cấu hình Docker Compose:
    \begin{lstlisting}[language=yaml]
services:
  postgres:
    image: postgres:15
    environment:
      POSTGRES_DB: analytics
    volumes: [pgdata:/var/lib/postgresql/data]
    healthcheck:
      test: ["CMD-SHELL", "pg_isready -U postgres"]
      interval: 5s
      timeout: 5s
      retries: 5
  worker:
    build: .
    depends_on:
      postgres:
        condition: service_healthy
    \end{lstlisting}
    \item \textbf{Bài 7}: Sử dụng cú pháp TaskFlow:
    \begin{lstlisting}[language=python]
@dag(schedule="@daily", start_date=datetime(2026, 1, 1), 
     default_args={"retries": 3, "retry_delay": timedelta(minutes=5)})
def etl_pipeline():
    @task
    def extract(): return [{"id": 1, "status": "COMPLETED", "amt": 100}]
    @task
    def transform(data): return [x for x in data if x["status"] != "CANCEL"]
    @task
    def load(clean): print(f"Loaded {len(clean)} rows.")
    load(transform(extract()))
etl_pipeline()
    \end{lstlisting}
    \item \textbf{Bài 8}: Cấu hình \texttt{trigger\_rule=TriggerRule.ALL\_DONE}. Mặc định Airflow là \texttt{all\_success} (nếu task trước fail thì task sau bị skip). Với \texttt{all\_done}, task thông báo sẽ luôn được thực thi sau khi tất cả upstream tasks đã hoàn tất chu kỳ chạy, bất kể thành công, thất bại hay bị hủy.
    \item \textbf{Bài 9}: dbt Model \texttt{stg\_orders.sql}:
    \begin{lstlisting}[language=sql]
select
    id as order_id,
    user_id as customer_id,
    cast(order_date as date) as order_date,
    amount_cents / 100.0 as order_amount_usd,
    status as order_status
from {{ source('raw_shop', 'orders') }}
where is_test = false
    \end{lstlisting}
    \item \textbf{Bài 10}: Trong \texttt{schema.yml}: Khai báo \texttt{columns: - name: customer\_id, tests: [unique, not\_null]}. Test tùy biến Singular (\texttt{tests/assert\_no\_negative\_orders.sql}):
    \begin{lstlisting}[language=sql]
SELECT order_id, order_total 
FROM {{ ref('stg_orders') }} 
WHERE order_total < 0
    \end{lstlisting}
    dbt sẽ báo lỗi (FAIL) nếu câu query này trả về bất kỳ dòng nào.
    \item \textbf{Bài 11}: Cấu hình dbt Incremental:
    \begin{lstlisting}[language=sql]
{{ config(materialized='incremental', unique_key='order_id') }}
select * from {{ ref('stg_orders') }}
{% if is_incremental() %}
    where updated_at > (select max(updated_at) from {{ this }})
{% endif %}
    \end{lstlisting}
    \item \textbf{Bài 12}: Ép Schema và ghi partitioned Parquet:
    \begin{lstlisting}[language=python]
schema = StructType([StructField("id", LongType()), StructField("val", DoubleType()),
                     StructField("year", IntegerType()), StructField("month", IntegerType())])
df = spark.read.schema(schema).csv("data/raw/*.csv")
df.write.partitionBy("year", "month").mode("overwrite").parquet("data/warehouse/fact/")
    \end{lstlisting}
    Tránh được thao tác quét toàn bộ dữ liệu (OOM crash) do inferSchema gây ra.
    \item \textbf{Bài 13}: Câu lệnh: \texttt{df\_joined = fact\_df.join(broadcast(dim\_store), "store\_id")}. Cơ chế: Spark sao chép toàn bộ bảng nhỏ (2MB) gửi đến bộ nhớ của tất cả các Worker Nodes, loại bỏ hoàn toàn giai đoạn Shuffle mạng của bảng lớn (500GB), chuyển Shuffle Hash Join thành Broadcast Hash Join với tốc độ nhanh gấp 10--50 lần.
    \item \textbf{Bài 14}: PySpark Window Functions:
    \begin{lstlisting}[language=python]
from pyspark.sql.window import Window
import pyspark.sql.functions as F

w_running = Window.partitionBy("employee_id").orderBy("date").rowsBetween(Window.unboundedPreceding, Window.currentRow)
df = df.withColumn("running_total", F.sum("sale_amount").over(w_running))

w_rank = Window.partitionBy("department").orderBy(F.desc("sale_amount"))
df_top3 = df.withColumn("rank", F.dense_rank().over(w_rank)).filter("rank <= 3")
    \end{lstlisting}
    \item \textbf{Bài 15}: Parquet chỉ là định dạng tệp lưu trữ dạng cột tĩnh không có transaction log. Delta Lake bổ sung tầng giao dịch ACID thông qua tệp nhật ký \texttt{\_delta\_log/}. Time travel: \texttt{SELECT * FROM delta\_table VERSION AS OF 3} hoặc \texttt{TIMESTAMP AS OF '2026-03-20'}. Lệnh \texttt{VACUUM} dùng để dọn dẹp vật lý các tệp dữ liệu rác cũ đã bị đánh dấu xóa vượt quá thời gian lưu giữ (retention period) để tiết kiệm dung lượng lưu trữ.
    \item \textbf{Bài 16}: Để đảm bảo thứ tự sự kiện của cùng 1 user, Producer bắt buộc phải đặt \texttt{partition\_key = user\_id}. Kafka đảm bảo toàn bộ message có cùng key sẽ rơi vào đúng 1 partition duy nhất (nơi thứ tự được bảo toàn 100\%). Consumer Lag là số lượng bản ghi đã được ghi vào Topic nhưng Consumer chưa kịp xử lý. Scale out: Tăng số lượng Consumer trong Consumer Group lên tối đa bằng số lượng Partitions (12 consumers).
    \item \textbf{Bài 17}: Ứng dụng FastAPI Serving chuẩn MLOps trong \texttt{code/ch12\_mlops/app.py}:
    \begin{lstlisting}[language=python]
from contextlib import asynccontextmanager
from fastapi import FastAPI
import joblib

ml_models = {}
@asynccontextmanager
async def lifespan(app: FastAPI):
    ml_models["churn"] = joblib.load("churn_model.joblib")
    yield
    ml_models.clear()

app = FastAPI(lifespan=lifespan)
@app.post("/predict")
def predict(payload: CustomerFeatures):
    prob = ml_models["churn"].predict_proba([[payload.total_spend, payload.tenure]])[0][1]
    return {"churn_prediction": int(prob >= 0.5), "churn_probability": float(prob)}
    \end{lstlisting}
    \item \textbf{Bài 18}: Kịch bản kiểm thử tải Locust:
    \begin{lstlisting}[language=python]
from locust import HttpUser, task, between

class MLUser(HttpUser):
    wait_time = between(0.1, 0.5)
    @task
    def predict_churn(self):
        payload = {"total_spend": 250.0, "tenure": 14}
        self.client.post("/predict", json=payload)
    \end{lstlisting}
    Chạy lệnh: \texttt{locust -f locustfile.py --headless -u 200 -r 20 --run-time 1m}. Mục tiêu: Đảm bảo P99 latency $< 50$ms và 0\% lỗi ở tải 500 RPS.
    \item \textbf{Bài 19}: Cấu hình GitHub Actions CI:
    \begin{lstlisting}[language=yaml]
name: Data Pipeline CI
on: [pull_request]
jobs:
  test:
    runs-on: ubuntu-latest
    steps:
      - uses: actions/checkout@v3
      - uses: actions/setup-python@v4
        with: { python-version: '3.11' }
      - run: pip install -r requirements-test.txt
      - run: flake8 . --max-line-length=100
      - run: pytest tests/ -v --cov=.
    \end{lstlisting}
    \item \textbf{Bài 20}: Data Drift: Phân phối của dữ liệu đầu vào $P(X)$ thay đổi qua thời gian. Concept Drift: Mối quan hệ giữa dữ liệu đầu vào và nhãn mục tiêu $P(Y \mid X)$ thay đổi. Chỉ số PSI: $PSI = \sum (\text{Actual}\% - \text{Baseline}\%) \times \ln(\text{Actual}\% / \text{Baseline}\%)$. Quy tắc MLOps: Nếu $PSI < 0.1$: Dữ liệu ổn định, không thay đổi đáng kể; Nếu $0.1 \le PSI < 0.25$: Bắt đầu có độ lệch nhẹ, cần theo dõi sát sao; Nếu $PSI \ge 0.25$: Độ lệch nghiêm trọng (Significant Drift), hệ thống MLOps phải tự động phát tín hiệu cảnh báo và kích hoạt pipeline huấn luyện lại mô hình (Model Retraining).
\end{itemize}
\end{dapanbox}
